\scanpage{36}
\uline{Proposition 2} Soit $f\in\operatorname{Aff}_B(G)$. \inserted{Supposons $G$ $A$-connexe.} Conditions équivalentes
\[
\begin{array}{ll}
\text{(i)}&f\in\operatorname{Aff}_A(G)\\
\text{(ii)}&f\mathbin{\lrcorner}U_A(G)\subset U_A(G)\\
\text{(iii)}&\langle f,U_A(G)\rangle\subset A
\end{array}
\]
\margin{\unsure{)}}

Évidemment
\diagram{\begin{tikzcd}[ampersand replacement=\&,column sep=small]
\text{(i)}\arrow[r,Rightarrow]\&\text{(ii)}\arrow[r,Rightarrow]\&\text{(iii)}
\end{tikzcd}}
[et
\diagram{\begin{tikzcd}[ampersand replacement=\&,column sep=small]
\deleted{\unsure{\text{(ii)}}}\inserted{\text{(iii)}}\arrow[r,Rightarrow]\&\text{(ii)}
\end{tikzcd}}
car \deleted{\unsure{soit}} $X\in U_A(G)$, pour prouver $f\mathbin{\lrcorner}X\in U_A(G)$ il suffit \inserted{(prop. 1)} de vérifier que $\langle g,f\mathbin{\lrcorner}X\rangle\in A$ pour $g\in\operatorname{Aff}_A(G)$, or
\[
\langle g,f\mathbin{\lrcorner}X\rangle=\langle gf,X\rangle=\langle f,g\mathbin{\lrcorner}X\rangle
\in\langle f,U_A(G)\rangle\subset A
\]
] Prouvons $(\mathrm{iii})\Rightarrow\mathrm{i})$ en utilisant que $G$ est $A$-connexe.

Soit $C_{\deleted{\unsure{A|B}}}=\operatorname{Aff}_A(G)$, considérons le diagramme
\diagram{\begin{tikzcd}[ampersand replacement=\&]
C\arrow[r]\arrow[d]\&C_B\arrow[d]\\
\widehat C\simeq\operatorname{Hom}_A(U,A)\arrow[r]\&\widehat C_B=\operatorname{Hom}_B(U_B,B)
\end{tikzcd}}
où $\widehat{\phantom C}$ désigne le \inserted{\unsure{séparé}} complété pour l'idéal d'augmentation.
\margin{35}
